\subsection{Веб-просмотрщик гауссиан}

В качестве демонстрации полученной модели выбрана реализация визуализации гауссиан с помощью WebGL (браузерный OpenGL) \cite{splat2024}. Так как исходный датасет содержал детальную разметку каждого органа и кости, данный репозиторий был доработан возможностью выделять определенные органы.

Для этого сначала требуется каждой гауссиане присвоить класс по принципу ближайшего соседа в размеченной воксельной сетке (листинг \ref{lst:segment-gaussians}).

Здесь же требуется рассмотреть особенный формат хранения гауссиан --- \texttt{.splat} по сравнению с \texttt{.ply}. Последний хранит гауссианы в том же формате, как они оптимизировались. Например, масштабы как логарифмы, кватернионы не обязательно единичной длины, непрозрачность как логит. Используемый для визуализации формат во первых с целью оптимизации переводит все эти характеристики в буквальные геометрические эквиваленты, а также из сферических гармоник собирает один, независимый от взгляда, цвет. В оригинальном репозитории есть реализация перевода \texttt{.ply} в \texttt{.splat} но она во-первых медленная (использует массивы python), а во-вторых не поддерживает метки классов, поэтому написана собственная реализация (листинг \ref{lst:ply-to-splat}).

Полный код форка представлен в приложении листингами \ref{lst:splat-html} (html) и \ref{lst:splat-js} (js).